\section{Introduction}
This chapter introduces the institutional context and project context for InlinePDF. It includes a profile of the college and a profile of the project, including tools, platform choices, and development setup.

\subsection{College Profile}
Bhagwan Mahavir College of Computer Application, Surat is a private institute established in 2008 by the Bhagwan Mahavir College Foundation, Surat. The institute is affiliated with Bhagwan Mahavir University and is governed by an executive board under the trust constitution and college by-laws.

The campus is located at New City Light Road, Bharthana-Vesu, Surat, in an open and green area. The institute includes classrooms, a documented library, a multipurpose hall, and modern computer laboratories. It serves approximately 1000 students.

\textbf{Vision:} The institute aims to prepare students for modern information technology environments by developing practical, skill-oriented capability and professional readiness.

\textbf{Mission:} The institute provides learning experiences in computer application through structured undergraduate education and practice-focused training for long-term career development.

\textbf{Quality objectives:} The institute aligns with the university mission of academic excellence through quality teaching, scholarship, and practical training. The department objective is to support both immediate career entry needs and long-term professional growth.

\subsection{Project Profile}
\textbf{Project title:} InlinePDF: A Local-First PDF Processing System for Privacy and Performance

\textbf{Project category:} Web-based application (local-first document utility)

\textbf{Problem domain:} PDF tools are widely used for merge, split, conversion, and organization tasks. Many public tools use a client-server model where uploaded files are processed on remote infrastructure. This model can introduce privacy concerns for sensitive documents and can add network-dependent delays.

\textbf{Project summary:} InlinePDF implements PDF operations directly in the browser. The system avoids server-side PDF processing and does not require user authentication, user accounts, or file upload pipelines for document operations. This approach aims to reduce unnecessary exposure of user files while maintaining practical speed for common tasks.

\textbf{Open-source position:} The project is open source. This supports implementation transparency, code auditability, and reproducibility for academic review.

\textbf{Development tools and environment:}
\begin{longtable}{p{0.30\textwidth}p{0.62\textwidth}}
\toprule
Item & Details \\
\midrule
Code editor & Visual Studio Code \\
Version control & Git and GitHub \\
Package manager & pnpm \\
Frontend framework & React with React Router and TypeScript \\
Build tooling & Vite \\
Styling and UI & Tailwind CSS and shadcn-based primitives \\
PDF libraries & pdf-lib and pdfjs-dist \\
Testing framework & Vitest and Testing Library \\
Deployment target & Cloudflare Workers \\
\bottomrule
\end{longtable}
